% ============================================================
%  INSTALACION Y REPRODUCCION
% ============================================================

\section{Instalación y Reproducción del Sistema}
\thispagestyle{fancy}

Este capítulo describe el procedimiento para poner en funcionamiento el servidor,
abrir el proyecto de Unity, generar la aplicación Android y reproducir una sesión
completa entre profesor y estudiantes. Se considera que el usuario ya dispone en su
computadora de las carpetas \texttt{Server} y \texttt{RealidadVirtual}. El teléfono
no resuelve los métodos numéricos, sino que recibe del servidor los resultados que
debe representar.

\subsection{Requisitos previos}

La instalación utiliza los elementos enumerados en la
Tabla~\ref{tab:requisitos-instalacion}.

\begin{table}[H]
\centering
\caption{Requisitos para reproducir el sistema.}
\label{tab:requisitos-instalacion}
{\small
\begin{tabular}{ll}
\toprule
\parbox[t]{4.0cm}{\textbf{Elemento}} & \parbox[t]{9.5cm}{\textbf{Requisito}} \\
\midrule
\parbox[t]{4.0cm}{Equipo servidor} & \parbox[t]{9.5cm}{Windows, Linux o macOS, conectado a la misma red local que los teléfonos.} \\
\parbox[t]{4.0cm}{Go} & \parbox[t]{9.5cm}{Versión \texttt{1.25.1}, de acuerdo con \texttt{Server/go.mod}.} \\
\parbox[t]{4.0cm}{Unity} & \parbox[t]{9.5cm}{Unity \texttt{6000.4.2f1} con Android Build Support, Android SDK, NDK y OpenJDK.} \\
\parbox[t]{4.0cm}{Teléfono} & \parbox[t]{9.5cm}{Android 8.0 (API 26) o superior, con pantalla táctil y sensor de orientación o giroscopio.} \\
\parbox[t]{4.0cm}{Red} & \parbox[t]{9.5cm}{Punto de acceso Wi-Fi sin aislamiento entre clientes; TCP 8080 y UDP 47777 habilitados.} \\
\parbox[t]{4.0cm}{Cable USB} & \parbox[t]{9.5cm}{Opcional; solo es necesario para instalar mediante ADB o depurar desde Unity.} \\
\bottomrule
\end{tabular}}
\end{table}

El teléfono cumple dos funciones: ejecuta el cliente Android y actúa como dispositivo
IHM. Se seleccionó porque constituye un recurso disponible para la totalidad de los
alumnos de la cátedra, lo que permite desarrollar la actividad con el curso completo
sin requerir la adquisición ni el préstamo de equipamiento especializado. La
interacción se realiza mediante la pantalla táctil y la IMU. No se requiere visor
Cardboard ni controlador externo. De forma
opcional puede utilizarse un Google Cardboard para activar la visualización
estereoscópica. El manejo integral mediante un mando externo se propone como trabajo
futuro y no forma parte de los requisitos para reproducir la versión entregada.

\newpage
\thispagestyle{fancy}
\subsection{Preparación y ejecución del servidor}

\subsubsection{Descarga de dependencias}

Desde una terminal ubicada en \texttt{Server} se ejecuta:

\begin{lstlisting}[basicstyle=\ttfamily\footnotesize, frame=single, numbers=none,
caption={Preparación del servidor Go.}]
go mod download
go build
\end{lstlisting}

La primera orden descarga Gin y las dependencias declaradas por el módulo. La segunda
comprueba que el servidor pueda compilarse en el equipo anfitrión.

\subsubsection{Inicio mediante el lanzador gráfico}

El modo recomendado se inicia desde la carpeta \texttt{Server} con:

\begin{lstlisting}[basicstyle=\ttfamily\footnotesize, frame=single, numbers=none,
caption={Inicio mediante el lanzador.}]
go run .
\end{lstlisting}

El programa abre una página local de configuración. Se conserva el puerto
\texttt{8080} y se pulsa el botón para iniciar. El lanzador muestra las direcciones IPv4
del equipo; estas sirven para la conexión manual si el descubrimiento automático no
está disponible.

Como alternativa, el servidor puede iniciarse directamente, sin lanzador:

\begin{lstlisting}[basicstyle=\ttfamily\footnotesize, frame=single, numbers=none,
caption={Inicio directo del servidor.}]
go run . -port 8080
\end{lstlisting}

Al arrancar se habilita la API HTTP en todas las interfaces de red y el respondedor de
descubrimiento en UDP 47777. En Windows, si aparece una consulta del firewall, se debe
permitir el acceso al menos para redes privadas. En una configuración manual se crean
reglas de entrada para TCP 8080 y UDP 47777.

\subsubsection{Verificación del panel docente}

En el equipo servidor se abre \texttt{http://localhost:8080}. El profesor ingresa el
código configurado en \texttt{main.go}; luego accede al panel, desde el cual puede ver
clientes, calcular soluciones y publicar o detener una sincronización. Para evitar
exponer credenciales en una distribución pública, se recomienda cambiar el código antes
de desplegar el sistema fuera de la red de laboratorio.

\subsection{Apertura y verificación del proyecto Unity}

\begin{enumerate}
    \item Abrir Unity Hub y comprobar que esté instalada la versión
    \texttt{6000.4.2f1} con Android Build Support y sus herramientas.
    \item Seleccionar \textit{Add project from disk} y elegir la carpeta
    \texttt{RealidadVirtual}.
    \item Esperar la importación de paquetes y la compilación de scripts.
    \item Abrir \textit{File $>$ Build Profiles} y seleccionar Android. Si todavía no
    está activo, utilizar \textit{Switch Platform}.
    \item Verificar que las escenas aparezcan habilitadas en este orden:\\
    {\small\texttt{MenuScene} $\rightarrow$ \texttt{HelloCardboard}
    $\rightarrow$ \texttt{HeatPlate}.}
\end{enumerate}

El proyecto está configurado para orientación automática, pantalla completa, API
mínima 26, permiso de Internet y tráfico HTTP en la red local. El identificador de
paquete puede cambiarse en \textit{Player Settings} si el APK debe coexistir con otra
compilación del mismo proyecto.

\subsection{Prueba previa dentro del editor}

Antes de generar el APK es conveniente comprobar el intercambio de datos
cliente--servidor:

\begin{enumerate}
    \item Iniciar el servidor y dejarlo activo.
    \item En Unity, abrir \texttt{MenuScene} y pulsar \textit{Play}.
    \item Desplegar la dirección manual e ingresar
    \texttt{http://127.0.0.1:8080}; el broadcast puede depender de la configuración
    del editor y no es necesario para esta prueba local.
    \item Completar nombre y legajo, y pulsar \textit{Conectar}.
    \item Confirmar en el panel docente que el alumno fue registrado y verificar que
    Unity cargue la escena de onda.
    \item Cambiar parámetros de onda y calor. El cliente debe enviar la petición,
    recibir todos los cuadros calculados y reproducirlos sin ejecutar el solver.
\end{enumerate}

La orientación por IMU no puede evaluarse de manera representativa con el ratón del
editor. Esa verificación se completa en el teléfono.

\subsection{Generación e instalación del APK}
\label{sec:compilar-apk}

\begin{enumerate}
    \item En \textit{Build Profiles}, confirmar Android y las tres escenas.
    \item En \textit{Edit $>$ Preferences $>$ External Tools}, utilizar las versiones
    de JDK, SDK y NDK instaladas con Unity para evitar incompatibilidades de Gradle.
    \item Pulsar \textit{Build}, elegir una carpeta de salida y guardar el archivo con
    un nombre identificable, por ejemplo \texttt{RV\_aplicacion.apk}.
    \item Copiar el APK al teléfono y permitir temporalmente la instalación desde esa
    fuente, o activar depuración USB e instalarlo con ADB:
\end{enumerate}

\begin{lstlisting}[basicstyle=\ttfamily\footnotesize, frame=single, numbers=none,
caption={Instalación opcional mediante ADB.}]
adb install -r RV_aplicacion.apk
\end{lstlisting}

La opción \texttt{-r} conserva los datos de una versión anterior. Para una prueba
completamente limpia se puede desinstalar previamente la aplicación desde Android.

\newpage
\thispagestyle{fancy}
\subsection{Reproducción de una sesión completa}

\begin{enumerate}
    \item Conectar la computadora del profesor y todos los teléfonos a la misma red
    Wi-Fi.
    \item Iniciar el servidor y comprobar el panel docente.
    \item Abrir la aplicación en cada teléfono, mantenerlo en horizontal, ingresar
    nombre y legajo, y pulsar \textit{Conectar}.
    \item Esperar la detección automática. Si no responde, desplegar el campo manual e
    ingresar la IPv4 mostrada por el lanzador, por ejemplo
    \texttt{http://192.168.1.20:8080}. En el teléfono nunca debe utilizarse
    \texttt{localhost}, porque esa dirección identifica al propio dispositivo.
    \item Confirmar desde el panel que todos los estudiantes aparecen registrados.
    \item En cada cliente, modificar parámetros, pausar o editar nodos y solicitar una
    nueva solución. Comprobar que la animación, la malla y el color se actualizan.
    \item Desde el panel docente, elegir onda o calor, establecer sus constantes y
    activar la sincronización. En aproximadamente 1,5 segundos, todos los clientes
    deben cambiar al modelo y versión publicados.
    \item Desactivar la sincronización y verificar que cada estudiante recupere el
    control local de sus parámetros.
\end{enumerate}

Durante la exploración, el usuario gira el teléfono para orientar la cámara y utiliza
la pantalla para operar paneles, mover la altura de la placa y seleccionar nodos. Esta
combinación materializa la IHM prevista para la actividad de Matemática Avanzada.
Si dispone de un Google Cardboard, puede pulsar \textit{MODO VR} antes de colocar el
teléfono en el visor. La cruz nativa de Cardboard permite regresar posteriormente a la
pantalla plana. En esta modalidad los paneles se ocultan, ya que su operación mediante
un mando Bluetooth constituye una ampliación futura.

\subsection{Diagnóstico de problemas frecuentes}

\begin{table}[H]
\centering
\caption{Comprobaciones ante fallas de instalación o conexión.}
\label{tab:diagnostico-instalacion}
{\small
\begin{tabular}{ll}
\toprule
\parbox[t]{4.2cm}{\textbf{Síntoma}} & \parbox[t]{9.3cm}{\textbf{Comprobación y corrección}} \\
\midrule
\parbox[t]{4.2cm}{El servidor no se detecta} & \parbox[t]{9.3cm}{Verificar la misma Wi-Fi, desactivar el aislamiento de clientes, habilitar UDP 47777 y probar la IPv4 manual.} \\
\parbox[t]{4.2cm}{No se conecta por IP} & \parbox[t]{9.3cm}{Usar la IP de la computadora, no \texttt{localhost}; comprobar TCP 8080 y que el servidor continúe activo.} \\
\parbox[t]{4.2cm}{HTTP 429} & \parbox[t]{9.3cm}{Esperar al menos cinco segundos. El servidor limita la frecuencia de cálculo y el cliente reintenta el último cambio.} \\
\parbox[t]{4.2cm}{Gradle falla al compilar} & \parbox[t]{9.3cm}{Seleccionar JDK, SDK y NDK incluidos con Unity; confirmar Android Build Support y espacio libre.} \\
\parbox[t]{4.2cm}{La escena no cambia} & \parbox[t]{9.3cm}{Revisar el orden de escenas y confirmar que el registro devolvió una respuesta satisfactoria.} \\
\parbox[t]{4.2cm}{La orientación no responde} & \parbox[t]{9.3cm}{Confirmar que el dispositivo posea sensor de rotación o giroscopio, mantenerlo horizontal y reiniciar la escena para recalibrar.} \\
\parbox[t]{4.2cm}{La sincronización no llega} & \parbox[t]{9.3cm}{Verificar \texttt{/sync/status}, la conectividad y que el profesor haya activado una nueva versión.} \\
\bottomrule
\end{tabular}}
\end{table}

\newpage
\thispagestyle{fancy}
\subsection{Lista de verificación final}

La reproducción se considera satisfactoria cuando se cumplen simultáneamente los
siguientes puntos:

\begin{itemize}
    \item el panel docente abre y el servidor responde en la red local;
    \item el teléfono descubre el servidor o se conecta mediante IPv4 manual;
    \item el registro del alumno aparece en el panel;
    \item onda y calor se calculan en Go y se reproducen en Unity;
    \item la pantalla táctil y la IMU permiten operar la IHM en el teléfono;
    \item los cambios docentes de modelo y parámetros se propagan a los clientes;
    \item la aplicación permite alternar entre vista plana a pantalla completa y un
    modo estereoscópico opcional para Google Cardboard;
    \item la cruz nativa de Cardboard permite regresar al modo plano, mientras que el
    control completo mediante mando queda documentado como trabajo futuro.
\end{itemize}
